Sometimes, one would like to classify \textbf{functions} rather than \textbf{languages} (remember, a computational task can be expressed by a function or a
language). This can be done generalizing some concepts we have previously introduced.

In particular, let $T: N \rightarrow N$ be the function used to detect how much time it's required for solving an instance of whatever problem, we state the
following definition.
\begin{definition}
    A function $f$ is in the class \textbf{FDTIME}(T(n)) if and only if \textbf{there is} a Turing Machine computing $f$ and running in time $n \rightarrow c 
    \cdot T(n)$ for some constant $c$.
\end{definition}
The complexity class \textbf{FP} is defined as below, very similar to the complexity class \textbf{P}:
$$ \textbf{FP} = \bigcup_{c \geq 1} \textbf{FDTIME}(n^c)$$
From the definition of the complexity class \textbf{FP} we can also describe other two peculiarities:
\begin{itemize}
    \renewcommand{\labelitemi}{-}
    \item For every language $L \in P$, the characteristic function $f$ of $L$ is trivially in \textbf{FP}.
    \item For certain classes of functions, there are canonical ways to turn a function $f$ into a language $L_f$. However, it is not true that $f \in \textbf{FP}$ implies $L_f \in \textbf{P}$.
\end{itemize}
